\documentclass[12pt]{article}
\usepackage[T1]{fontenc}
\usepackage[utf8]{inputenc}
\usepackage[brazil]{babel}
\usepackage[margin=1in]{geometry}
\usepackage{ragged2e}
\usepackage[normalem]{ulem}
\setlength{\parindent}{0pt}
\setlength{\parskip}{0.8\baselineskip}
\begin{document}
\justifying
\noindent Verifica-se que a questão objeto da controvérsia jurídica suscitada no recurso manejado pela parte recorrente esteve afetada no representativo da controvérsia - \textbf{Tema n. }\textbf{331} da TNU - (PU no processo n. PEDILEF 5008761-19.2020.4.04.7102/RS),  afetado em 14/06/2023 foi julgado (trânsito em julgado em 13/09/2024), por meio do qual foi elucidada a seguinte questão:

\noindent \textit{"Determinar se, no caso de movimentações bancárias fraudulentas realizadas por terceiro, mediante uso de cartão magnético e senha pessoal do correntista, pode caracterizar falha de segurança do banco, apta a afastar a excludente de responsabilidade do art. 14, § 3º, inciso II, do Código de Defesa do Consumidor, a ausência de verificação da autenticidade das referidas movimentações, quando atípicas e/ou suspeitas em relação ao perfil do correntista."}

\noindent Restou firmada a seguinte tese:

\noindent \textit{"1. O uso indevido de cartão de débito ou crédito por terceiro, mediante fraude, constitui, em regra, fortuito interno para os fins da Súmula 479/STJ, salvo se comprovada culpa exclusiva do consumidor ou de terceiro (art. 14, § 3º, inciso II, do Código de Defesa do Consumidor). 2. Em princípio, a realização de operação com o uso de cartão e senha descaracteriza a responsabilidade do banco por configurar quebra do dever contratual de cuidado do cliente. 3. Todavia, não se configura a excludente de responsabilidade se, independentemente de prévia comunicação da ocorrência pelo titular do cartão, (i) as circunstâncias em que as operações foram realizadas e o perfil do consumidor revelarem fortes indícios de fraude detectáveis pelo banco; ou (ii) não restar claramente demonstrado o descumprimento consciente, pelo consumidor, do dever contratual de cuidado no uso do cartão, seja em razão do grau de sofisticação dos meios de engenharia social empregados pelos fraudadores, seja pela condição de hipervulnerabilidade da vítima." (Tema 331 TNU).}

\noindent Com efeito, depreende-se da leitura do voto condutor do acórdão que o mesmo se encontra em conformidade com o decidido no  tema pela TNU.

\noindent A proposito, consta do acordao hostilizado: (...) Dever de Informação, DIREITO DO CONSUMIDOR, Crédito Direto ao Consumidor - CDC, Bancários, Contratos de Consumo, DIREITO DO CONSUMIDOR

\noindent Nessas condições, o conteúdo do Acórdão recorrido é consentâneo com o entendimento da TNU, firmado em julgamento de recurso representativo de controvérsia, o que atrai a incidência da regra prevista pelo art. 14, III, \textit{b,} da Resolução CJF n. 586/2019:

\noindent \textit{\textbf{Art. 14.}}\textit{ Decorrido o prazo para contrarrazões, os autos serão conclusos ao magistrado responsável pelo exame preliminar de admissibilidade, que deverá, de forma sucessiva: (...) }\textit{\textbf{III }}\textit{- negar seguimento a pedido de uniformização de interpretação de lei federal interposto contra acórdão que esteja em conformidade com entendimento consolidado: }\textit{\textbf{b)}}\textit{ em }\uline{\textit{\textbf{recurso representativo de controvérsia pela Turma Nacional de Uniformização}}}\textit{ ou em pedido de uniformização de interpretação de lei dirigido ao Superior Tribunal de Justiça;(...).}

\noindent Posto isso, com arrimo no \uline{\textbf{art. 14, III, }}\uline{\textit{\textbf{b}}}\textit{,} da Resolução CJF n. 586/2019 c/c art. 1.030, I, do CPC, \textbf{NEGO SEGUIMENTO }\textbf{a}o\textbf{ PU}\textbf{.}

\noindent Intimem-se. Aguarde-se o decurso do prazo recursal. Após, certifique-se o trânsito em julgado e devolva-se ao juízo de origem.
\end{document}
